\chapter{Hook-to-timing cheat sheet}
\label{app:hook-timing}
\index{hook timing}\index{after-init-hook}\index{emacs-startup-hook}\index{window-setup-hook}\index{find-file-hook}

\noindent
Table~\ref{tab:hook-timing} summarises each hook's firing point and whether it contributes to \texttt{emacs-init-time}.

\begin{sidewaystable}[p]
\centering
\caption{Hook timing reference}
\label{tab:hook-timing}
\begin{tabularx}{\linewidth}{@{}lXll@{}}
\toprule
Hook & When & In \texttt{emacs-init-time}? & In \texttt{-{}-batch}? \\
\midrule
(top-level forms in init.el) & During init.el eval & Yes & Yes \\
\texttt{after-init-hook} & After init.el finishes & Yes & Yes \\
\texttt{emacs-startup-hook} & After \texttt{after-init-hook}, after CLI args & Yes & Yes \\
\texttt{window-setup-hook} & After frame is visible & No & \textbf{No} \\
\texttt{run-with-idle-timer} (any) & After N seconds of idle & No & \textbf{No} \\
\texttt{find-file-hook} & Per file open & N/A & Yes \\
\texttt{after-change-major-mode-hook} & Per mode activation & N/A & Yes \\
\bottomrule
\end{tabularx}
\end{sidewaystable}
